% Part: second-order-logic
% Chapter: syntax-and-semantics
% Section: expressive-power

\documentclass[../../../include/open-logic-section]{subfiles}

\begin{document}

\olfileid{sol}{syn}{exp}
\olsection{వ్యక్తీకరణశక్తి}

\begin{explain}
ద్వితీయ-స్థాయి చరాలపై పరిమాణీకరణ, మొదటిస్థాయి తర్కంతో పోలిస్తే భాష
వ్యక్తీకరణశక్తిని అపారంగా పెంచుతుంది. కొన్ని ధర్మాలు గల ప్రమేయాలు లేదా
సంబంధాలు ఉన్నాయని చెప్పడానికి ద్వితీయ-స్థాయి అస్తిత్వ పరిమాణీకరణ వీలు
కల్పిస్తుంది. మొదటిస్థాయి తర్కంలో అలా చెప్పడానికి, ఆ ప్రయోజనం కోసం
ఒక తార్కికేతర సంకేతాన్ని (అంటే \tetoken{ఒక ప్రమేయ సంకేతం}{!!a{function}} లేదా
\tetoken{ఒక విధేయ సంకేతాన్ని}{!!{predicate}}) నిర్దేశించడం మాత్రమే మార్గం. \tetoken{వ్యక్తి క్షేత్రం}{!!{domain}}
యొక్క అన్ని ఉపసమితులు, దానిపై అన్ని సంబంధాలు, లేదా \tetoken{వ్యక్తి క్షేత్రం}{!!{domain}}
నుంచి దానికే అన్ని ప్రమేయాలు ఏదో ఒక ధర్మాన్ని కలిగి ఉన్నాయని
చెప్పడానికి ద్వితీయ-స్థాయి సార్వత్రిక పరిమాణీకరణ వీలు కల్పిస్తుంది.
మొదటిస్థాయి తర్కంలో, భాషలోని ఏదో ఒక తార్కికేతర సంకేతానికి నిర్దేశించిన
ఉపసమితులు, సంబంధాలు లేదా ప్రమేయాలు ఒక ధర్మాన్ని కలిగి ఉన్నాయని మాత్రమే
చెప్పగలం. ఒక ధర్మం గల ఉపసమితులు, సంబంధాలు, ప్రమేయాలు ఉన్నాయని లేదా
వాటన్నింటికీ ఆ ధర్మం ఉందని చెప్పేటప్పుడు, ఆ ధర్మాన్ని నిర్దేశించడంలోనూ
ద్వితీయ-స్థాయి పరిమాణీకరణను ఉపయోగించవచ్చు. దీనివల్ల మొదటిస్థాయి
తర్కంలో నిర్వచించలేని సంబంధాలను నిర్వచించవచ్చు; మొదటిస్థాయి తర్కంలో
వ్యక్తం చేయలేని వ్యక్తి క్షేత్ర ధర్మాలను వ్యక్తం చేయవచ్చు.
\end{explain}

\begin{defn}
$\Struct{M}$ అనేది భాష~$\Lang{L}$కు \tetoken{ఒక నిర్మాణం}{!!a{structure}} అయితే,
స్వేచ్ఛగా $x$,~$y$ చరాలు మాత్రమే గల ఏదో \tetoken{సూత్రం}{!!{formula}}~$!A_R(x, y)$
ఉండి, $s(x) = a$, $s(y) = b$లకు $R(a, b)$ అవడం (అంటే
$\tuple{a,b} \in R$ అవడం), $\Sat{M}{!A_R(x, y)}[s]$ అవడం సమానార్థకం
అయితే, సంబంధం~$R \subseteq \Domain{M}^2$ను~$\Lang{L}$లో
\emph{నిర్వచించదగినది} అంటాం.
\end{defn}

\begin{ex}
మొదటిస్థాయి తర్కంలో తాదాత్మ్య సంబంధం~$\Id{\Domain{M}}$ను (అంటే
$\Setabs{\tuple{a,a}}{a \in \Domain{M}}$ను) $\eq[x][y]$ అనే
సూత్రంతో నిర్వచించవచ్చు. ద్వితీయ-స్థాయి తర్కంలో, ఈ సంబంధాన్ని
\emph{$\eq$ లేకుండానే} నిర్వచించవచ్చు. $a$,~$b$లు~$\Domain{M}$లో
ఒకే \tetoken{మూలకం}{!!{element}} అయితే, అవి~$\Domain{M}$ యొక్క అవే ఉపసమితుల్లో
\tetoken{మూలకాలుగా}{!!{element}s} ఉంటాయి (ఎందుకంటే సమితులను వాటి \tetoken{మూలకాలే}{!!{element}s}
నిర్ణయిస్తాయి). విపర్యయంగా, $a$,~$b$లు వేర్వేరైతే అవి అవే
  ఉపసమితుల్లో \tetoken{మూలకాలుగా}{!!{element}s} ఉండవు: ఉదాహరణకు, $a \neq b$ అయితే
$a \in \{a\}$ కానీ $b \notin \{a\}$. అందువల్ల
``$\Domain{M}$ యొక్క అవే ఉపసమితుల్లో \tetoken{మూలకాలుగా}{!!{element}s} ఉండటం'' అనే
సంబంధం $a = b$ అయినప్పుడు, అప్పుడు మాత్రమే $a$,~$b$లకు వర్తిస్తుంది.
$\Domain{M}$ యొక్క అన్ని ఉపసమితులపై పరిమాణీకరించగలం కాబట్టి, ఈ
సంబంధాన్ని ద్వితీయ-స్థాయి తర్కంలో వ్యక్తం చేయవచ్చు. కాబట్టి కింది
\tetoken{సూత్రం}{!!{formula}} $\Id{\Domain{M}}$ను నిర్వచిస్తుంది:
\[
\lforall[X][(X(x) \liff X(y))]
\]
\end{ex}

\begin{prob}
$\lforall[X][(X(x) \lif X(y))]$ కూడా (గమనించండి:
$\liff$ కాదు,~$\lif$!) $\Id{\Domain{M}}$ను నిర్వచిస్తుందని చూపండి.
\end{prob}

\begin{ex}
$R$ ఒక ద్విస్థాన \tetoken{విధేయ సంకేతం}{!!{predicate}} అయితే, $\Assign{R}{M}$ అనేది
~$\Domain{M}$పై ఒక ద్విస్థాన సంబంధం. కొంత గందరగోళం కలిగించే
విధంగా, \tetoken{విధేయ సంకేతం}{!!{predicate}}~$R$కు, సంబంధం~$\Assign{R}{M}$కూ $R$నే
ఉపయోగిస్తాం. $R$ యొక్క \emph{సంక్రమణ సంవృతి}~$R^*$ అంటే, ఏవో
$c_1$, \dots, $c_k$లకు $R(a,c_1)$, $R(c_1, c_2)$, \dots,
$R(c_k,b)$ అయినప్పుడు, అప్పుడు మాత్రమే $a$,~$b$ల మధ్య వర్తించే
సంబంధం. ఇందులో $k = 0$ సందర్భం కూడా ఉంటుంది; అంటే $R(a,b)$ అయితే
$R^*(a,b)$ కూడా. కాబట్టి $R \subseteq R^*$. నిజానికి, $R$ను కలిగి
ఉండి సంక్రమణమైన సంబంధాల్లో $R^*$ అతి చిన్నది. $X$ అనేది $R$ను కలిగి
ఉన్న సంక్రమణ సంబంధమని ద్వితీయ-స్థాయి తర్కంలో ఇలా చెప్పవచ్చు:
\begin{multline*}
  !B_R(X) \ident \lforall[x][\lforall[y][(R(x,y) \lif X(x, y))]] \land {}\\
\lforall[x][\lforall[y][\lforall[z][((X(x,y) \land X(y,z)) \lif X(x,
      z))]]].
\end{multline*}
మొదటి సంయుక్తాంశం $R \subseteq X$ అని, రెండవది $X$ సంక్రమణమని
చెబుతాయి.

  అది అటువంటి సంబంధాల్లో అతి చిన్నదని చెప్పడమంటే, $R$ను కలిగి ఉండి
సంక్రమణమైన ప్రతి సంబంధంలోనూ $X$ కూడా ఇమిడి ఉందని చెప్పడమే. అందువల్ల
$R$ యొక్క సంక్రమణ సంవృతిని కింది \tetoken{సూత్రంతో}{!!{formula}} నిర్వచించవచ్చు:
\[
R^*(X) \ident !B_R(X) \land \lforall[Y][(!B_R(Y) \lif
  \lforall[x][\lforall[y][(X(x, y) \lif Y(x,y))]])].
\]
$\Sat{M}{R^*(X)}[s]$ అవుతుంది అప్పుడే, అప్పుడు మాత్రమే $s(X) = R^*$.
$R$ యొక్క సంక్రమణ సంవృతిని మొదటిస్థాయి తర్కంలో వ్యక్తం చేయలేం.
\end{ex}

\end{document}
